\section*{الفصل \ref{ch:b2:pt-paths} --- القياس الحراري الضغطي ومسارات الضغط--الحرارة--الزمن}

\begin{solution}{exo:b2:pt-paths:1}
(أ) A نطاق الكلوريت، وB نطاق البيوتيت، وC نطاق الغارنت، وD نطاق الستوروليت، وE نطاق الكيانيت،
وF نطاق السيليمانيت. (ب) خمسة: خطوط تساوي الدرجة للبيوتيت والغارنت والستوروليت والكيانيت والسيليمانيت،
كلٌّ منها بين عيّنتين متتاليتين، والدرجة ترتفع من A إلى F. (ج) خط تساوي الدرجة
لا يقول إلا إن تفاعلًا جرى؛ أمّا التدرّج الحقلي فيحتاج إلى ضغوط ودرجات حرارة،
تأتي من تفاعلات معايَرة أو من القياس الحراري الضغطي لمجموعة الذروة في كل
صخر.
\end{solution}

\begin{solution}{exo:b2:pt-paths:2}
(أ) $(X_{\text{Mg}}/X_{\text{Fe}})_{\text{الغارنت}} = 0.10/0.80 = 0.125$،
$(X_{\text{Mg}}/X_{\text{Fe}})_{\text{البيوتيت}} = 0.45/0.50 = 0.900$؛ $K_D = 0.1389$، $\ln K_D =
-1.9741$. (ب) $3R = 24.942$؛ والمقام $20.19 + 24.942\times1.9741 = 69.428$؛ عند
\qty{5}{kbar}، $T = 54185/69.428 = \qty{780.5}{K} = \qty{507}{\celsius}$؛ وعند
\qty{7}{kbar}، $T = (52630 + 2177)/69.428 = \qty{789.4}{K} = \qty{516}{\celsius}$.
(ج) كيلوباران لا يغيّرانها إلا بمقدار \qty{9}{K}: فهو ميزان حرارة.
\end{solution}

\begin{solution}{exo:b2:pt-paths:3}
(أ) $z = P/(\bar\rho g) = 6.0\times10^8/(2800\times9.81) = \qty{21.8}{km}$.
(ب) $(580 - 10)/21.8 = \qty{26}{\celsius\per\km}$. (ج) عند \qty{580}{\celsius} يقع الخط
عند $3.67 + 0.8\times2.10 = \qty{5.35}{kbar}$؛ وعند \qty{6.0}{kbar} يكون الصخر في
جهة الضغط المرتفع: الكيانيت.
\end{solution}

\begin{solution}{exo:b2:pt-paths:4}
(أ) المتفاعلات $2\times140.693 + 3\times18.015 = \qty{335.431}{g}$؛ والنواتج $277.112 +
58.320 = \qty{335.432}{g}$: موزونة في حدود التقريب. (ب) تصير $2\times43.65 =
\qty{87.30}{cm^3}$ مقدار $107.50 + 24.63 = \qty{132.13}{cm^3}$: $+\qty{51.4}{\percent}$.
(ج) الأوليفين $281.386/87.30 = \qty{3.223}{g.cm^{-3}}$؛ والنواتج $335.432/132.13 =
\qty{2.539}{g.cm^{-3}}$. (د) ينتفخ الصخر بنصف حجمه ويصير أخفّ؛
ويخزن أيضًا الماء في معادنه.
\end{solution}

\begin{solution}{exo:b2:pt-paths:5}
(أ) المتوسطان $-1.206$ عند \qty{1072.15}{K} و$-1.782$ عند \qty{823.15}{K}. ومع $1/T =
9.3271\times10^{-4}$ و$1.21485\times10^{-3}$، يكون $a = (-1.206 + 1.782)/(9.3271\times10^{-4}
- 1.21485\times10^{-3}) = -2041.5$ كلفن و$b = -1.206 + 2041.5/1072.15 = 0.6982$.
(ب) $\Delta_r H + 2070\times0.311 = 2041.5\times24.942 = \qty{50920}{J}$، ومنه $\Delta_r H =
\qty{50276}{J}$؛ $\Delta_r S = 0.6982\times24.942 = \qty{17.41}{J.K^{-1}}$. (ج) عند
\qty{5}{kbar}: من أجل $K_D = 0.18$، 588.1 مقابل \qty{587.5}{\celsius}؛ ومن أجل $0.10$، 419
مقابل \qty{425}{\celsius}؛ ومن أجل $0.30$، 819 مقابل \qty{806}{\celsius}. فضمن مجال
التجارب تتوافق المعايرتان، مع أن ثوابتهما تختلف بمقدار
\qty{2.4}{kJ} و\qty{2.8}{J.K^{-1}}؛ إذ تتعاوض أخطاء $\Delta_r H$ و$\Delta_r S$
هناك وتكبر خارجه. فلا ينبغي استعمال معايرة بعيدًا
عن ظروف تجاربها.
\end{solution}

\begin{solution}{exo:b2:pt-paths:6}
(أ) التبادل: $(113.12 + 154.30) - (115.32 + 149.65) = \qty{+2.45}{cm^3} =
\qty{0.245}{J.bar^{-1}}$، مقابل 0.311. وGASP: $125.28 + 2\times44.15 + 22.69 -
3\times100.79 = \qty{-66.10}{cm^3} = \qty{-6.61}{J.bar^{-1}}$، مقابل $-6.21$، وهي قيمة
المعايرة عند \qty{800}{\celsius} و\qty{10}{kbar}، حيث يغيّر التمدّد الحراري
والانضغاط الأحجام. وفي الحالتين يبلغ تغيّر GASP عشرين إلى ثلاثين ضعف
تغيّر التبادل. (ب) $\partial T/\partial P = 0.311/62.96 = \qty{4.9e-3}{K/bar}$:
\qty{4.9}{K} لكل كيلوبار. (ج) ميل الخط $22.80 + 8.314\times3\ln 0.20/6.21
= 22.80 - 6.46 = \qty{16.3}{bar/K}$؛ و\qty{17}{K} تُزيح الضغط بمقدار
\qty{280}{bar}.
\end{solution}

\begin{solution}{exo:b2:pt-paths:7}
(أ) $P_0 = 22.80\times600 - 1093 = \qty{12587}{bar}$؛ $3\ln(0.08/0.30) = -3.9653$؛ $RT =
8.314\times873.15 = \qty{7259.4}{J}$؛ $P = 12587 + 7259.4\times(-3.9653)/6.21 = 12587 -
4635 = \qty{7952}{bar}$. (ب) $3\ln(0.08/0.35) = -4.4277$، $P = 12587 - 5176 =
\qty{7411}{bar}$: أقل بمقدار \qty{0.54}{kbar}، فيجب تحليل البلاجيوكليز بعناية.
(ج) $7952\times10^5/(2800\times9.81) = \qty{29.0}{km}$.
\end{solution}

\begin{solution}{exo:b2:pt-paths:8}
(أ) $z = P/(\bar\rho g)$: \qty{118.9}{km} و\qty{34.0}{km} و\qty{15.3}{km}.
(ب) $(118.9 - 34.0)/2.2 = \qty{38.6}{km/Myr} = \qty{3.9}{cm.yr^{-1}}$ و$(34.0 -
15.3)/1.1 = \qty{17.0}{km/Myr} = \qty{1.7}{cm.yr^{-1}}$، أعلى قليلًا من القيمتين المنشورتين
3.4 و1.6: إذ يعطي تحويلهم للضغط إلى عمق فروق أعماق
أصغر قليلًا. (ج) $\sigma = \sqrt{0.9^2 + 0.9^2} = \qty{1.3}{Myr}$، ومنه $2.2 \pm 1.3$
مليون سنة ومعدّل بين $84.9/3.47 = 24$ و$84.9/0.93 = \qty{91}{km/Myr}$،
أي من 2.4 إلى \qty{9.1}{cm.yr^{-1}}: سريع في كل الأحوال. (د) حتى أدنى معدّل يبلغ نحو خمسين ضعف
قيمة النموذج \qty{0.5}{mm.yr^{-1}}: فالحتّ وحده لا يستطيع ذلك.
\end{solution}

\begin{solution}{exo:b2:pt-paths:9}
(أ) $\partial T/\partial\ln K_D = 3RT/(\Delta_r S - 3R\ln K_D) = 24.942\times860.6/62.96 =
\qty{341}{K}$؛ $\partial T/\partial P = 0.311/62.96 = \qty{4.9}{K}$ لكل كيلوبار.
(ب) $\sigma_T = \sqrt{(341\times0.05)^2 + 4.9^2} = \sqrt{17.0^2 + 4.9^2} =
\qty{17.7}{K}$. (ج) الخطأ التحليلي على $K_D$ هو الغالب؛ وهو وحده ثلثا
عدم اليقين البالغ \qty{25}{\celsius} لأفضل المعايرات، الذي يشمل أيضًا أخطاء
المعايرة نفسها.
\end{solution}

\begin{solution}{exo:b2:pt-paths:10}
(أ) الغارنت $0.10/0.75 = 0.1333$، والبيوتيت $0.42/0.52 = 0.8077$، $K_D = 0.1650$، $\ln K_D =
-1.8018$، والمقام $20.19 + 24.942\times1.8018 = 65.130$؛ $3\ln(0.07/0.25) = -3.8189$.
عند \qty{6000}{bar}: $T = 54496/65.130 = \qty{836.7}{K} = \qty{563.6}{\celsius}$، ثم $P
= 22.80\times563.6 - 1093 + 8.314\times836.7\times(-3.8189)/6.21 = \qty{7479}{bar}$.
وعند \qty{7479}{bar}: $T = \qty{570.6}{\celsius}$، $P = \qty{7603}{bar}$. وعند
\qty{7603}{bar}: $T = \qty{571.2}{\celsius}$، $P = \qty{7614}{bar}$؛ ومرة أخرى، $T =
\qty{571.3}{\celsius}$. فقد بلغ الصخر التوازن عند نحو \qty{571}{\celsius}
و\qty{7.6}{kbar}. (ب) عند \qty{571}{\celsius} يقع خط الكيانيت $=$ السيليمانيت عند $3.67
+ 0.71\times2.10 = \qty{5.16}{kbar}$: الكيانيت. (ج) ميزان الضغط معايَر من أجل
الكيانيت؛ ومع السيليمانيت يختلف التفاعل، ويختلف خطه، لأن متعدّدي الشكل لا
يملكان طاقة جيبس نفسها.
\end{solution}

\begin{solution}{exo:b2:pt-paths:11}
(أ) $\tau = L^2/(\pi^2\kappa)$: من أجل \qty{35}{km}، $1.225\times10^9/(9.870\times10^{-6}) =
\qty{1.24e14}{s} = \qty{3.9}{Myr}$؛ ومن أجل \qty{70}{km}، أربعة أضعاف ذلك،
\qty{15.7}{Myr}. (ب) سبعون مليون سنة تساوي من أربعة إلى ثمانية عشر ضعف هذين الزمنين: فالحرارة
تنتقل عبر القشرة المثخنة أسرع بكثير مما يزيلها الحتّ. (ج) يبلغ الصخر
أكبر عمق له دفعة واحدة، عند التكديس؛ وتصل الحرارة على مدى عشرات ملايين السنين
التالية، بينما يرفعه الحتّ ببطء: فتظلّ درجة الحرارة
ترتفع مدةً أثناء انخفاض الضغط. (د) حتٌّ في سبعة ملايين سنة كان
سيرفع الصخر قبل أن يسخّنه التوصيل: ذروة حرارية أدنى، تُبلغ
مبكرًا، عند ضغط أعلى، ومسار بارد عائد إلى السطح.
\end{solution}

\begin{solution}{exo:b2:pt-paths:12}
(أ) تزن القشرة $35000\times2800\times9.81 = \qty{0.961}{GPa}$؛ وتحتاج القيمة الباقية
\qty{1.839}{GPa} إلى $1.839\times10^9/(3300\times9.81) = \qty{56.8}{km}$ من الوشاح: فالعمق
\qty{91.8}{km}، كما تقول الدراسة الأصلية: \qty{90}{km} أو أكثر. (ب) $(4.0 -
0.961)\times10^9/(3300\times9.81) = \qty{93.9}{km}$، والعمق \qty{128.9}{km}.
(ج) $700/91.8 = \qty{7.6}{\celsius\per\km}$: نسبة مرتفعة جدًّا للضغط إلى
درجة الحرارة، أي متتالية نطاقات الاندساس. (د) $(91.8 - 35)/34 = \qty{1.7}{Myr}$.
\end{solution}

\begin{solution}{pb:b2:pt-paths:1}
\noindent\textbf{الجزء الأول.}

\textbf{1.} نواة الغارنت $0.080/0.760 = 0.1053$، والبيوتيت $0.40/0.58 = 0.6897$؛ $K_D =
0.1527$، $\ln K_D = -1.8793$.

\textbf{2.} الحافة $0.120/0.720 = 0.1667$، وبيوتيت الأرضية $0.46/0.53 = 0.8679$؛ $K_D =
0.1921$، $\ln K_D = -1.6497$.

\textbf{3.} $3R = 24.942$؛ والبسط $52630 + 0.311\times6000 = 54496$؛ والمقام $20.19
+ 24.942\times1.8793 = 67.064$؛ $T = \qty{812.6}{K} = \qty{539.5}{\celsius}$.

\textbf{4.} المقام $20.19 + 24.942\times1.6497 = 61.337$؛ $T = 54496/61.337 =
\qty{888.5}{K} = \qty{615.3}{\celsius}$.

\textbf{5.} $0.311\times1000/67.064 = \qty{4.6}{K}$ للنواة و$0.311\times1000/61.337
= \qty{5.1}{K}$ للحافة.

\textbf{6.} ينطبق ميزان الحرارة على معادن كانت في توازن معًا. فالبيوتيت
المحتبس في النواة عُزل عن بقية الصخر حين نمت النواة؛
أمّا بيوتيت الأرضية فآخر ما تبادل كان مع الحافة. وقرن النواة بالأرضية يخلط
لحظتين.

\noindent\textbf{الجزء الثاني.}

\textbf{7.} النواة $22.80\times539.5 - 1093 = \qty{11208}{bar}$؛ والحافة $22.80\times615.3 -
1093 = \qty{12936}{bar}$.

\textbf{8.} $3\ln(0.050/0.25) = -4.8283$؛ $RT = 8.314\times812.65 = \qty{6756.4}{J}$؛ $P =
11208 - 6756.4\times4.8283/6.21 = 11208 - 5253 = \qty{5955}{bar}$.

\textbf{9.} $3\ln(0.060/0.28) = -4.6213$؛ $RT = 8.314\times888.45 = \qty{7386.6}{J}$؛ $P
= 12936 - 5497 = \qty{7439}{bar}$.

\textbf{10.} عند \qty{7439}{bar}: $T = (52630 + 2314)/61.337 = \qty{895.8}{K} =
\qty{622.6}{\celsius}$؛ وعند درجة الحرارة هذه $P = 13102 - 5542 = \qty{7560}{bar}$؛ ثم $T
= \qty{623.2}{\celsius}$. فقد نمت الحافة عند نحو \qty{623}{\celsius} و\qty{7.56}{kbar}.

\textbf{11.} عند \qty{5954}{bar}، $T = \qty{539.2}{\celsius}$ و$P = \qty{5950}{bar}$: فقد
نمت النواة عند نحو \qty{539}{\celsius} و\qty{5.95}{kbar}. وكان التخمين الأول،
\qty{6}{kbar}، قريبًا أصلًا، ولا يتحرّك ميزان الحرارة إلا \qty{4.6}{K} لكل كيلوبار.

\textbf{12.} النواة $5950\times10^5/(2800\times9.81) = \qty{21.7}{km}$؛ والحافة
$7560\times10^5/27468 = \qty{27.5}{km}$.

\textbf{13.} عند \qty{539}{\celsius} يقع الخط عند $3.67 + 0.39\times2.10 =
\qty{4.49}{kbar}$، تحت \qty{5.95}{kbar}: الكيانيت، وهو متّسق مع الكيانيت المحتبس
في النواة. وعند \qty{623}{\celsius} يقع عند $5.77 + 0.23\times2.09 = \qty{6.25}{kbar}$،
تحت \qty{7.56}{kbar}: الكيانيت مرة أخرى.

\noindent\textbf{الجزء الثالث.}

\textbf{14.} $+\qty{84}{\celsius}$ و$+\qty{1.6}{kbar}$ (أعمق بنحو \qty{6}{km}): كان
الصخر يُدفن ويُسخَّن، على \omterm{def:b2:pt-paths:paths}{مساره التقدّمي}.

\textbf{15.} عند \qty{623}{\celsius} أو أكثر، تحت \qty{6.25}{kbar}: بعد نموّ الحافة
فقد الصخر \qty{1.3}{kbar} على الأقل، أي نحو \qty{5}{km} من الحمل العلوي، مع بقائه على
الأقل بالسخونة نفسها. فقد استُخرج وهو لا يزال ساخنًا.

\textbf{16.} الكلوريت معدن مائي منخفض الدرجة: يسجّل \omterm{def:b2:pt-paths:paths}{المسار التراجعي}
المتأخر، أي التبريد على عمق ضحل، مع الماء.

\textbf{17.} دفن وتسخين حتى الحافة، ثم انخفاض الضغط دون تبريد، ثم التبريد:
إذا رُسم مع درجة الحرارة إلى اليمين والضغط إلى الأعلى، ينعطف المسار باتجاه عقارب الساعة.
وجاءت ذروة الضغط أولًا، عند الحافة؛ والذروة الحرارية في الوقت نفسه أو بعده، أثناء
انخفاض الضغط.

\textbf{18.} نعم: فالنموذج أيضًا يبلغ ذروته الضغطية أولًا وذروته الحرارية،
\qty{593}{\celsius} عند \qty{7.8}{kbar}، أثناء الصعود، قريبًا من ظروف حافة
الشيست. فالإثخان متبوعًا بحتّ بطيء يفسّر مسار الشيست.

\noindent\textbf{الجزء الرابع.}

\textbf{19.} $(623 - 405)/(412.0 - 398.0) = 218/14.0 = \qty{15.6}{\celsius}$ لكل مليون
سنة.

\textbf{20.} برد الصخر أسرع بنحو مرة ونصف مما افتُرض. والمعدن
الذي يبرد أسرع يملك وقتًا أقل ليفقد أرغونه، فينغلق عند درجة حرارة أعلى:
فدرجة حرارة الانغلاق الحقيقية تقع فوق \qty{405}{\celsius} بقليل، والانغلاق
أعمق بقليل.

\textbf{21.} $(405 - 10)/30 = \qty{13.2}{km}$.

\textbf{22.} $27.5 - 13.2 = \qty{14.3}{km}$ (وبالأعماق غير المقرّبة، 27.52 و13.17،
\qty{14.35}{km}).

\textbf{23.} $\sqrt{1.5^2 + 1.0^2} = \qty{1.8}{Myr}$: $14.0 \pm \qty{1.8}{Myr}$، ومنه معدّل
بين $14.35/15.8 = 0.91$ و$14.35/12.2 = \qty{1.18}{km/Myr}$.

\textbf{24.} $14.35/14.0 = \qty{1.03}{km}$ لكل مليون سنة، أي
\textbf{\qty[text-series-to-math=true]{1.0}{mm.yr^{-1}}}، أبطأ بنحو ثلاثين مرة من
المرحلة الأولى لاستخراج أعمق صخور الألب، ونحو ضعف معدّل الحتّ
في نموذج الإثخان.
\end{solution}
